\documentclass{article}
\usepackage{PRIMEarxiv} % Core package for the PrimeAI Template
\usepackage{amsmath, amssymb, amsthm, bm, dcolumn} % Add amsthm here for the proof environment
\usepackage[numbers,sort&compress]{natbib} % Natbib for citations
\usepackage{graphicx} % For high-quality images
\usepackage{multicol} % Multi-column figures/tables
\usepackage{hyperref} % Hyperlinks
\usepackage{listings} % Code listings
\usepackage{authblk} % For structured affiliations
% Define theorem style (optional)
\theoremstyle{plain}
\newtheorem{theorem}{Theorem}


% Custom commands for primes and qubit encoding
\newcommand{\primeQubit}[1]{|p_{#1}\rangle}
\newcommand{\primeGate}[1]{U_{p_{#1}}}

\usepackage{wrapfig}
\usepackage[pscoord]{eso-pic}
\usepackage[fulladjust]{marginnote}
\reversemarginpar

% Typesetting improvements without footnote patching
\usepackage[protrusion=true, expansion=true, tracking=false]{microtype}
\microtypecontext{spacing=nonfrench}

% Line numbers
\usepackage[right]{lineno}

% Text layout - adjust as needed
\raggedright
\setlength{\parindent}{0.5cm}
\textwidth 5.25in 
\textheight 8.75in

% Set double spacing
\usepackage{setspace} 
\doublespacing

% Adjust width for specific content
\usepackage{changepage}

% Adjust caption style
\usepackage[aboveskip=1pt,labelfont=bf,labelsep=period,singlelinecheck=off]{caption}

% Remove brackets from references
\makeatletter
\renewcommand{\@biblabel}[1]{\quad#1.}
\makeatother

% Header, footer, and page numbers
\usepackage{lastpage,fancyhdr}
\pagestyle{fancy}
\fancyhf{}
\fancyhead[L]{Preprint - PrimeAI Enhanced Template}
\fancyfoot[C]{\scriptsize Multiplicity Theory © 2024 Ryan Van Gelder - Citizen Gardens\\ Licensed Under MIT and CC BY-NC-SA 4.0.}
\fancyfoot[R]{Page \thepage\ of \pageref{LastPage}}
\renewcommand{\footrule}{\hrule height 2pt \vspace{2mm}}

\begin{document}

\title{Integrating Peter Shor's Contributions with Multiplicity Theory}
\author{Ryan O. Van Gelder \\ Citizen Gardens - The Foundation of Multiplicity}
\date{\today}
\maketitle

\begin{abstract}
Multiplicity theory presents a novel mathematical framework rooted in prime number theory and multiset-theoretic principles. By integrating Peter Shor's contributions, particularly his algorithm for quantum factorization, this paper demonstrates how the Matrix Compute Paradigm (MCP) enhances computational efficiency through prime encoding, modified quantum gates, and entanglement optimization. This integration yields novel methods for addressing challenges in quantum speedup, error resilience, and scalable architectures.
\end{abstract}

\section{Introduction}
Peter Shor's algorithm for integer factorization revolutionized quantum computation by demonstrating exponential speedup over classical algorithms \cite{shor1994algorithms, shor1997polynomial}. Within the MCP framework, Shor’s insights are further adapted by utilizing prime encoding and multiplicative structures to enhance computational efficiency and error resilience.

\section{Prime Encoding in Quantum States}
Each qubit is represented as a prime number $p_i$, forming a prime-encoded quantum state:
\begin{equation}
|\psi\rangle = |p_i\rangle
\end{equation}
This encoding leverages the unique multiplicative properties of primes to simplify modular operations and reduce factorization errors. Parameters $i_k$ are mapped to primes $p_k$ for computational efficiency:
\begin{equation}
f(i_k) = p_k, \quad \text{where } i_k \in I.
\end{equation}

\section{Modified Quantum Fourier Transform (QFT)}
The QFT is adapted for prime-encoded states. For a quantum state $|p_i\rangle$, the transformation becomes:
\begin{equation}
\text{QFT}: |p_i\rangle \to \frac{1}{\sqrt{P}} \sum_{k=0}^{P-1} e^{2\pi i p_i k / P} |p_k\rangle,
\end{equation}
where $P$ is a prime modulus. The modified QFT supports phase estimation based on prime periods, enhancing Shor's algorithm's periodicity detection.

\section{Prime-Based Quantum Gates}
Prime-controlled gates, such as phase gates, maintain the prime-encoded structure:
\begin{equation}
U_{p_i} = \sum_{j=0}^{P-1} e^{2\pi i p_i j / P} |j\rangle \langle j|.
\end{equation}
Modular exponentiation in this framework:
\begin{equation}
U_{\text{mod}} |x\rangle = |(x \cdot a) \mod N\rangle,
\end{equation}
optimizes computational overhead by utilizing prime states.

\section{Entanglement and Superposition}
Prime-encoded states in superposition:
\begin{equation}
\Psi(t) = \sum_{i=1}^N \alpha_i \Psi_i e^{i\theta_i(t)},
\end{equation}
enable parallel computations with reduced interference. Entanglement is represented as:
\begin{equation}
E(|p_i\rangle |p_j\rangle) = \frac{1}{\sqrt{2}} \left(|p_i\rangle |p_j\rangle + |p_j\rangle |p_i\rangle\right).
\end{equation}

\section{Error Resilience with Feedback Loops}
Recursive feedback dynamically adjusts quantum states:
\begin{equation}
F(p_i) = p_i \cdot (1 + \delta),
\end{equation}
where $\delta$ accounts for error margins. Tensor networks manage entanglement and ensure scalable, error-resilient computations.

\section{Conclusion}
Shor’s algorithm within the MCP, augmented by prime encoding and modular operations, offers breakthroughs in quantum speedup, error resilience, and scalable architecture design. This synergy underscores the transformative potential of integrating Shor's contributions into MCP \cite{nielsen2010quantum, preskill2018quantum}.

\bibliographystyle{plain}
\bibliography{references}
\end{document}
